\documentclass[11pt]{article}
\usepackage[margin=1in]{geometry}
\usepackage{amsmath}
\usepackage{booktabs}

\title{Step 4 Validation Protocol: Independent Temperature and Field Trends}
\author{}
\date{}

\begin{document}
\maketitle

\section{Purpose of the test}
Step 3 showed that a dynamic dimerisation/phonon coordinate can produce a
nonlinear response that is distinct from a purely scalar bright-dark mixing.
Step 4 tests the trend-level interpretation needed before assigning a measured
feature to spin-Peierls physics.  The control is:
\begin{equation}
\delta(T,B) \rightarrow 0
\quad\mathrm{above\ the\ spin\!-\!Peierls\ transition},
\end{equation}
while the local spin correlation proxy remains finite,
\begin{equation}
C_1(T,B) \ne 0
\quad\mathrm{in\ both\ uniform\ and\ dimerised\ regimes}.
\end{equation}
The second condition is essential: if $C_1$ were forced to zero above the
transition, the test would no longer distinguish a lattice order parameter from
a uniform spin-correlation channel.

\section{Physical model}
The Hamiltonian is the Step 3 model,
\begin{equation}
H_{\mathrm{mix}} =
V_{BD}^{\mathrm{static}} L_{BD}
+ g_Q^{\mathrm{eff}} Q L_{BD},
\end{equation}
with
\begin{equation}
V_{BD}^{\mathrm{static}} =
V_0 + \lambda_\delta \delta(T,B) + \lambda_C C_1(T,B).
\end{equation}
The dynamic lattice coupling is scaled by the spin-Peierls order parameter,
\begin{equation}
g_Q^{\mathrm{eff}} = g_{Q0}
\frac{\delta(T,B)}{\delta(T_{\mathrm{ref}},B_{\mathrm{ref}})}.
\end{equation}
Thus the dynamic lattice channel turns off when the dimerisation turns off.

\section{Trend definitions}
The spin-Peierls transition temperature is represented phenomenologically by
\begin{equation}
T_{\mathrm{SP}}(B) =
\max\left[0, T_{\mathrm{SP},0}(1-\alpha_B B^2)\right].
\end{equation}
The dimerisation is
\begin{equation}
\delta(T,B)=
\begin{cases}
\delta_0 \left(1-T/T_{\mathrm{SP}}(B)\right)^\beta,
& T<T_{\mathrm{SP}}(B),\\
0, & T\ge T_{\mathrm{SP}}(B).
\end{cases}
\end{equation}
The spin-correlation proxy is not an order parameter.  It is kept finite above
the transition and varies smoothly with temperature and field.

\section{Scenarios}
\begin{center}
\begin{tabular}{llll}
\toprule
Scenario & Regime & Lattice channel & Purpose \\
\midrule
\texttt{lowT\_lattice\_dynamic} & low $T$, low $B$ & on & spin-Peierls reference \\
\texttt{highT\_lattice\_off\_C1\_finite} & high $T$ & off & uniform chain with finite $C_1$ \\
\texttt{highB\_lattice\_suppressed} & low $T$, high $B$ & off & field-suppressed order \\
\texttt{highT\_C1\_static} & high $T$ & off & persistent scalar $C_1$ signal \\
\bottomrule
\end{tabular}
\end{center}

\section{Acceptance criteria}
A lattice-controlled feature should be large in the low-temperature
spin-Peierls case and suppressed in both the high-temperature and
field-suppressed controls.  A feature that remains large in the
\texttt{highT\_C1\_static} case is compatible with a spin-correlation channel,
but it should not be interpreted uniquely as evidence for the dimerised phase.

\section{Requested outputs}
For each scenario, save:
\begin{itemize}
\item the rephasing real spectrum PDF;
\item the pathway-resolved complex response arrays;
\item the rephasing, unrephasing, and 1Q total components;
\item text and JSON summaries containing $T$, $B$, $T_{\mathrm{SP}}$,
$\delta$, $C_1$, $g_Q^{\mathrm{eff}}$, and $V_{BD}^{\mathrm{static}}$.
\end{itemize}
The analysis should compare finite-window integrals, full-spectrum differences,
and trend metrics across the four scenarios.

\section{Numerical results}
The four Step 4 scenarios were generated at $\eta=0.001$ with a
$120\times120$ frequency grid, $n_{\mathrm{bosons}}=3$, and
$\omega_Q=0.035$ eV.  The finite-window half-width used for the aggregate
metrics is $0.025$ eV.

\begin{center}
\begin{tabular}{lrrrrr}
\toprule
Scenario & $\delta$ & $C_1$ & $g_Q^{\mathrm{eff}}$ &
$V_{BD}^{\mathrm{static}}$ & $A_{\mathrm{cross}}$ \\
\midrule
Low-$T$ lattice dynamic & 0.007071 & -0.378667 & 0.0100 & 0.0200 & 12.554241 \\
High-$T$ lattice off & 0 & -0.329471 & 0 & 0.0100 & 10.919127 \\
High-$B$ lattice suppressed & 0 & -0.346199 & 0 & 0.0100 & 10.919127 \\
High-$T$ $C_1$ static & 0 & -0.329471 & 0 & 0.0200 & 16.266623 \\
\bottomrule
\end{tabular}
\end{center}

The low-temperature lattice reference has finite $\delta$ and finite
$g_Q^{\mathrm{eff}}$.  The high-temperature and high-field lattice-off controls
both have $\delta=0$ and $g_Q^{\mathrm{eff}}=0$, while $C_1$ remains finite.
Because these two controls also have the same $V_{BD}^{\mathrm{static}}=0.01$,
their spectra are identical in this minimal model.  The full-spectrum relative
maximum difference between them is exactly zero at numerical precision.

The low-temperature lattice spectrum differs strongly from the high-temperature
lattice-off control.  The full-spectrum relative maximum difference is
$6.76\times10^{-1}$, and the main cross-window area changes from $12.554241$ to
$10.919127$.  This confirms that the implemented lattice channel follows the
spin-Peierls trend.

The high-temperature $C_1$ static control is the key interpretive warning.  It
keeps $\delta=0$ and $g_Q^{\mathrm{eff}}=0$, but uses finite $C_1$ to restore
$V_{BD}^{\mathrm{static}}=0.02$.  Its main cross-window area is
$16.266623$, larger than both lattice-off controls.  The full-spectrum relative
maximum difference between the high-temperature lattice-off control and this
$C_1$ static control is $5.98\times10^{-1}$.

\section{Validation outcome}
The trend-control validation passes.  The model now contains two independent
trend channels: a dimerisation/lattice channel that vanishes above the
spin-Peierls transition or under field suppression, and a spin-correlation
proxy that remains finite in the uniform regime.

The conclusion is therefore sharper than in Step 3.  A feature that disappears
when $\delta(T,B)$ and $g_Q^{\mathrm{eff}}$ vanish is compatible with a
spin-Peierls/lattice-controlled mechanism.  A feature that persists in the
high-temperature $C_1$ static control, where $\delta=0$, cannot be assigned
uniquely to the dimerised phase.  It may instead reflect a scalar local
spin-correlation contribution to bright-dark mixing.

\section{Conclusions and next step}
Step 4 provides the minimal trend-level filter needed before connecting the
model to CuGeO3 data.  The lattice channel should be judged by whether it tracks
$\delta(T,B)$, not merely by whether the signal is present at low temperature.
The finite-$C_1$ control shows that a nonlinear feature can remain strong even
when the spin-Peierls order parameter is zero.

The next useful calculation is to turn this discrete control into a scan:
\begin{itemize}
\item sweep $T$ across $T_{\mathrm{SP}}$ at fixed low field and track selected
window integrals;
\item sweep $B$ at fixed low temperature to check field suppression of
$\delta(T,B)$;
\item compare lattice-dynamic and $C_1$-static channels using the same plotted
observables rather than only full spectra;
\item only after that, add a momentum-sector or spinon-inspired treatment if
the trend-level signatures remain distinct.
\end{itemize}

\end{document}
